\section{Discussion}
\label{sec:discussion}

\subsection{Against the hypothesis}
The hypothesis has three parts. That restricting scope increases trust is supported in direction by the questionnaire (65 percent of 374), by the 30 prospective users (26 expressed trust, 10 fully) and by the simulation (none disagreed), and the campaign shows the boundary mechanisms hold under hostile use; the support is conditional on the system behaving exactly as promised. That it increases willingness to adopt is partly supported, resting on stated intention after a pitch and a short trial, and would not be supported under the strict reading of the simulation. That it does so without meaningfully reducing usefulness has not been tested, because no task measure exists yet; what can be said is that the limits were not what users objected to, and that reliability, not restriction, currently bounds the experience. The hypothesis as a whole is encouraged and not confirmed.

\subsection{Structure is necessary and not sufficient}
The prototype shows that the restriction can be built as structure rather than policy, cheaply, and that multi-step plans, tidying, cited answers, mail, search, drafting, voice and remote control can all run through three enforcement points. It also shows that a passing automated suite did not guarantee a working interface: 71 defects were found while the suite was green, and several of them, a stale banner, an unlabelled unsupported claim, a microphone that stayed live, would damage trust as surely as a boundary violation. The guarantee a user experiences is the composition of the enforcement core and everything wired to it, and the second must be tested through the interface, under adversarial use, and not inferred from the first.

\subsection{Threats to validity}
The evidence is of four kinds with different weaknesses, which is why they are reported separately and not merged. The deterministic tests cannot exercise a real model. The interface campaign used agents that cannot hear and cannot feel, on a machine lacking the embedding model. The 30 conversations were conducted by the team that built and presented the system, which risks a favourable response, and were not recorded with individual scores. The simulation is synthetic. The questionnaire sample is self-selected, student-heavy and drawn from one institution's networks, and it rated a description rather than a demonstration. None of these is the controlled study.

\subsection{Planned controlled study}
\label{sec:study}
The second phase first corrects the recorded defects, with a regression test for each, and re-runs the campaign; a study of a prototype that is still changing would measure the changes and not the idea. The study then compares participants completing representative multi-file tasks with and without the system, recording completion time and success, and instruments trust specifically against the bounded-scope claim so that the hypothesis is tested rather than assumed, alongside a subjective workload instrument and a usability questionnaire. The baseline is the same assistant with the boundary and gates relaxed, so that the restriction is the manipulated variable. Three further measurements address the inherited limitations: approval fatigue over repeated sessions, since a single session cannot exhibit it; synthesis faithfulness and citation accuracy with a real model, reported as a known limitation where they fall short; and voice against a four-second latency target with workload and trust alongside task time. Classification criteria for the hypothesis and each supporting claim are fixed before the data are collected.
